% !TEX root = ../main.tex

\documentclass[../main.tex]{subfiles}

\graphicspath{{\subfix{../images/}}}

\begin{document}
	
\subsection{Общая структура системы визуализации}

На основе опыта, полученного при реализации модуля чтения, была
разработана полноценная система \textbf{Pixhawk IMU Visualizer}~---
многопоточное приложение на C++. Два параллельных потока взаимодействуют
через потокобезопасный разделяемый буфер (\texttt{SharedBuffer}):
 
\textit{Главный поток} (\texttt{pixhawk\_imu\_reader.cpp}) устанавливает
соединение с Pixhawk по UART, непрерывно парсит входящий байтовый поток
и при получении пакета \texttt{RAW\_IMU} помещает данные в очередь, а также
дублирует их в CSV-файл.
 
\textit{Поток визуализации} (\texttt{imu\_visualizer.cpp}) извлекает пакеты
из очереди, запускает алгоритм инерциальной навигации и обновляет 3D-сцену на
каждом кадре.


\subsection{Структура файлов проекта}

\begin{table}[H]
\centering
\caption{Файлы проекта}
\label{tab:files}
\begin{tabularx}{\textwidth}{|l|X|}
\hline
\textbf{Файл} & \textbf{Назначение} \\
\hline
\texttt{imu\_shared.h}            & Общие структуры данных: \texttt{RawIMU},
                                    \texttt{NavState}, \texttt{SharedBuffer},
                                    константы калибровки \\
\hline
\texttt{imu\_visualizer.h}        & Объявление функции \texttt{run\_visualizer()} \\
\hline
\texttt{imu\_visualizer.cpp}      & Навигационный алгоритм, GLSL-шейдеры,
                                    геометрия модели и OpenGL-рендер \\
\hline
\texttt{pixhawk\_imu\_reader.cpp} & MAVLink-ридер, запуск потока визуализации \\
\hline
\texttt{raw\_data.cpp}            & Модуль первичного считывания \\
\hline
\texttt{CMakeLists.txt}           & Система сборки CMake \\
\hline
\texttt{c\_library\_v2/}          & Заголовочная библиотека MAVLink~C~\cite{mavlink_docs} \\
\hline
\end{tabularx}
\end{table}


\subsection{Разделяемые структуры данных $(\texttt{imu\_shared.h})$}

Вся общая информация инкапсулирована в трёх структурах. Структура
\texttt{RawIMU} хранит одно необработанное измерение:
 
\begin{lstlisting}[caption={Структуры RawIMU и NavState}]
struct RawIMU {
    double  timestamp;
    int16_t xacc,  yacc,  zacc;   // acceleration (LSB)
    int16_t xgyro, ygyro, zgyro;  // gyroscope (LSB)
    int16_t xmag,  ymag,  zmag;   // magnetometer (LSB)
};

struct NavState {
    double qw=1, qx=0, qy=0, qz=0; // quaternion orientation
    double wx=0, wy=0, wz=0;        // angular velocity (rad/s)
    double vx=0, vy=0, vz=0;        // velocity (m/s)
    double px=0, py=0, pz=0;        // position (m)
    double roll=0, pitch=0, yaw=0;  // Euler angles (deg)
    double ax=0, ay=0, az=0;       // acceleration (m/s^2)
};
\end{lstlisting}
 
Структура \texttt{SharedBuffer} объединяет всё межпоточное взаимодействие.
Каждое поле защищено отдельным мьютексом, что позволяет потокам работать
параллельно без глобальной блокировки. Флаг \texttt{running} объявлен
как \texttt{std::atomic<bool>}:
 
\begin{lstlisting}[caption={Потокобезопасный разделяемый буфер}]
struct SharedBuffer {
    std::mutex           mtx;
    std::deque<RawIMU>   queue;
    static const int     MAX_QUEUE = 500;
    std::atomic<bool>    running{true};

    std::mutex           nav_mtx;
    NavState             nav;

    std::mutex           trail_mtx;
    std::deque<std::array<float,3>> trail;
    static const int     MAX_TRAIL = 2000;

    void push(const RawIMU& r) {
        std::lock_guard<std::mutex> lk(mtx);
        if ((int)queue.size() < MAX_QUEUE) queue.push_back(r);
    }

    bool pop(RawIMU& out) {
        std::lock_guard<std::mutex> lk(mtx);
        if (queue.empty()) return false;
        out = queue.front();
        queue.pop_front();
        return true;
    }
};
\end{lstlisting}

	
\end{document}
